% Mechanistic Views Audit — standalone \input{} file
% Requires: booktabs, array, multirow, xcolor, cleveref, float packages in parent document

\section{Mechanistic Views Classification}
\label{app:mechviews}

Mechanistic Views \citep{tower2026mechviews} classifies mechanistic claims across five
axes: \emph{ontology} (what kind of entity), \emph{identity} (when two
descriptions pick out the same thing), \emph{evidence} (what measurements
warrant the claim), \emph{formalism} (what math expresses it), and
\emph{target} (what phenomenon is explained). Nine views partition the space
of mechanistic claims; a given paper may employ several views, either
deliberately or through implicit \emph{view slides}---unannounced shifts
between ontological commitments. This appendix classifies the paper's central
claims and flags the most consequential slides.

\paragraph{Dominant view and key slides.}
The paper's primary ontology is the \textbf{subspace view}: the causal
variable is a point on $\Gr(k,d)$, and the atlas systematically tests whether
each operation's causal variable lives on this manifold. The formalism
($\Gr(k,d)$, geodesic distance, principal angles) is fully native to this
view. However, three consequential view slides occur:

\emph{Subspace $\to$ structural view} (\S\ref{sec:gauge_freedom},
\S\ref{sec:svd_vs_das}). The subspace degeneracy result---10 seeds produce
orthogonal subspaces that all support $\IIA \geq 0.94$---undermines subspace
\emph{identity} as an axis of comparison. The paper responds by shifting to a
structural view: what matters is not which subspace, but the gauge-invariant
information flow through the weight matrices (SVD predicts $\IIA = 0.999$
without optimization). This slide is acknowledged in the text (``the identity
of a causal subspace is not a meaningful property'') but not framed as a view
change, which it is: the relevant entity shifts from a specific point on
$\Gr(k,d)$ to an orbit under the gauge group.

\emph{Subspace $\to$ process view} (\S\ref{sec:mechanism}). The mechanistic
analysis shifts from asking ``what is the causal variable?'' (subspace view)
to ``how does it emerge?'' (process view). The AGOP double-peak, spectral
compression, and Fourier lag analyses all operate in the process view with
dynamical-systems formalism. This is a natural and productive shift, but it
means the mechanistic section does not further validate the subspace claims---it
explains a different aspect of the same phenomenon.

\emph{Contrastive view as implicit methodology.} The memorization artifact
analysis (\S\ref{sec:memorization}) and NL-DAS vacuity analysis
(\S\ref{sec:nldas_vacuous}) are inherently contrastive: they gain evidential
force by contrasting genuine causal structure with artifactual alternatives.
The 4-level hierarchy itself is a contrastive framework (each level defines
what the level below fails to exclude). This contrastive methodology is
pervasive but never named as a view.

\paragraph{The productive tension.}
The paper's most important finding from a Mechanistic Views perspective is that its
own primary ontology is self-undermining: the subspace view defines the atlas
and motivates the Grassmannian formalism, but the degeneracy result shows that
subspace identity is gauge-dependent. This is not a flaw---it is the paper's
strongest argument for the structural view (weight-space SVD, gauge-invariant
composites like $W_Q W_K^\top$) as the correct level of description for
causal geometry. The 4-level evaluation hierarchy maps directly onto the
evidence axis of Mechanistic Views, providing a graduated standard for when
patching-based evidence (levels 1--2) must be supplemented by distributional
(level 3) and structural (level 4) measurements.

\vspace{1em}
\begin{table}[H]
\centering
\caption{Mechanistic Views audit of this paper's central claims.
Dominant view in bold; $\to$ indicates a view slide within the claim.}
\label{tab:mechviews}
\small
\setlength{\tabcolsep}{4.5pt}
\renewcommand{\arraystretch}{1.25}
\begin{tabular}{@{}p{4.8cm}p{1.8cm}p{1.8cm}p{2.0cm}p{1.6cm}p{1.4cm}@{}}
\toprule
\textbf{Claim} & \textbf{Ontology} & \textbf{Identity} & \textbf{Evidence} & \textbf{Formalism} & \textbf{Target} \\
\midrule
Three-class atlas partition
  & \textbf{Subspace} & Subsp.\ overlap & Patch, Probe & $\Gr(k,d)$ & Behavior \\
Grokking governs the partition
  & Process & Func.\ equiv. & Train.\ dyn. & Dyn.\ sys. & Algorithm \\
DAS subspace degeneracy (10 seeds)
  & Sub$\to$Struct & Gauge inv. & Patch, Weight & $\Gr(k,d)$ & Repr. \\
SVD predicts causal geometry
  & \textbf{Structural} & Gauge inv. & Weight & $\Gr(k,d)$ & Repr. \\
Memorization artifacts ($\IIA$ w/o structure)
  & \textbf{Contrastive} & Func.\ equiv. & Patch & $\Gr$, Func. & Behavior \\
NL-DAS vacuity (diversity ratio)
  & Contrastive & Func.\ equiv. & Patch & Functions & Behavior \\
Struct.\ pi-SAE recovers NL variables
  & Subspace & Func.\ equiv. & Patch, Probe & Functions & Algorithm \\
DMF+ReLU groks 18$\times$ faster
  & Process & Func.\ equiv. & Train.\ dyn. & Dyn.\ sys. & Capacity \\
Fourier structure lags generalization
  & Process & --- & TD, Weight & Dyn.\ sys. & Algorithm \\
4-level evaluation hierarchy
  & --- & --- & Patch, Probe & Functions & Behavior \\
\bottomrule
\end{tabular}
\renewcommand{\arraystretch}{1.0}
\end{table}
